\subsection{Seven-curve $I_D$--$V_D$ output-characteristic family (22 October 2026)}
The assignment's supplied \texttt{SDevice\_IdVd.txt} command file was
run using the same completed SProcess geometry, with
\texttt{Grid = "@tdr@"} preprocessing to \texttt{n5\_fps.tdr}.
A separate Workbench copy (\texttt{SOI\_nMOSFET\_Copy}) preserved
the original dual-bias $I_D$--$V_G$ simulations. The output-characteristic
sequence initializes the gate at zero bias, saves seven gate-bias states,
and loads each state to sweep the drain from $V_{DS}=0$ to $1.2$ V.
The gate biases are $V_{GS}=0,0.2,0.4,0.6,0.8,1.0,1.2$ V.
The 0-V initial gate setting corrects the original source script's
inconsistent initial value of $-0.5$ V before its first saved state.
The SDevice solver reported successful completion on 22 October 2026
at 02:36:46 SGT in 647.15 s (10 min 47 s), with all seven
\texttt{idvd\_vgXX\_n3\_des.plt} files generated.

The original seven-curve Sentaurus Visual CSV
(\texttt{Id\_vs\_Vd\_00\_12.csv}) contains 31 equally spaced drain-bias
values for each of the seven curves (increments of $0.04$ V).
The seven pairs of independent X/Y columns were consolidated into
\texttt{data/idvd\_baseline\_seven\_gate\_biases.csv}, with one shared
drain-voltage column and seven distinct drain-current columns.
No interpolation or smoothing was applied.

\begin{figure}[H]\centering
\begin{tikzpicture}
\begin{axis}[
width=0.94\linewidth,height=7cm,
xlabel={$V_{DS}$ (V)},ylabel={$I_D$ (exported current units)},
xmin=0,xmax=1.2,ymin=0,ymax=0.00062,
grid=both,legend pos=north west,
legend style={font=\small,draw=none},
scaled y ticks=false,
ytick={0,0.0001,0.0002,0.0003,0.0004,0.0005,0.0006},
yticklabels={$0$,$1\cdot10^{-4}$,$2\cdot10^{-4}$,$3\cdot10^{-4}$,$4\cdot10^{-4}$,$5\cdot10^{-4}$,$6\cdot10^{-4}$}
]
\addplot+[thick,no marks]
table[col sep=comma,x=VDS_V,y=Id_VGS_00_A_exported]
{data/idvd_baseline_seven_gate_biases.csv};
\addlegendentry{$V_{GS}=0.0$ V}
\addplot+[thick,no marks]
table[col sep=comma,x=VDS_V,y=Id_VGS_02_A_exported]
{data/idvd_baseline_seven_gate_biases.csv};
\addlegendentry{$V_{GS}=0.2$ V}
\addplot+[thick,no marks]
table[col sep=comma,x=VDS_V,y=Id_VGS_04_A_exported]
{data/idvd_baseline_seven_gate_biases.csv};
\addlegendentry{$V_{GS}=0.4$ V}
\addplot+[thick,no marks]
table[col sep=comma,x=VDS_V,y=Id_VGS_06_A_exported]
{data/idvd_baseline_seven_gate_biases.csv};
\addlegendentry{$V_{GS}=0.6$ V}
\addplot+[thick,no marks]
table[col sep=comma,x=VDS_V,y=Id_VGS_08_A_exported]
{data/idvd_baseline_seven_gate_biases.csv};
\addlegendentry{$V_{GS}=0.8$ V}
\addplot+[thick,no marks]
table[col sep=comma,x=VDS_V,y=Id_VGS_10_A_exported]
{data/idvd_baseline_seven_gate_biases.csv};
\addlegendentry{$V_{GS}=1.0$ V}
\addplot+[thick,no marks]
table[col sep=comma,x=VDS_V,y=Id_VGS_12_A_exported]
{data/idvd_baseline_seven_gate_biases.csv};
\addlegendentry{$V_{GS}=1.2$ V}
\end{axis}
\end{tikzpicture}
\caption{Simulated SOI nMOSFET output characteristics: drain
\texttt{TotalCurrent} versus drain \texttt{OuterVoltage} for seven gate
biases. All curves share the same current axis. The nearly zero initial
$V_{DS}=0$ currents reflect the simulation's numerical zero-current floor.}
\label{fig:idvd_seven}
\end{figure}

\begin{table}[H]\centering
\caption{Direct values at $V_{DS}=1.2$ V, extracted from the supplied
Sentaurus Visual CSV. Current magnitudes are in raw exported units;
the corresponding mA values assume the exported units represent A.}
\begin{tabular}{@{}lll@{}}\toprule
$V_{GS}$ (V) & $I_D$ (exported units) & $I_D$ (mA, assuming A)\\\midrule
0.0 & $3.4856612\times10^{-10}$ & $3.486\times10^{-7}$\\
0.2 & $4.0002410\times10^{-7}$ & 0.0004000\\
0.4 & $2.8921286\times10^{-5}$ & 0.028921\\
0.6 & $1.2513808\times10^{-4}$ & 0.125138\\
0.8 & $2.6207930\times10^{-4}$ & 0.262079\\
1.0 & $4.1793077\times10^{-4}$ & 0.417931\\
1.2 & $5.8128743\times10^{-4}$ & 0.581287\\
\bottomrule\end{tabular}
\end{table}

Increasing $V_{GS}$ raises the drain current and steepens the initial
low-$V_{DS}$ region. At sufficiently high drain bias, the current
becomes much less sensitive to $V_{DS}$ (the saturation-like region).
Its remaining positive slope indicates finite output conductance,
consistent with effects such as channel-length modulation; the curve
alone does not uniquely isolate the physical mechanism.
For instance, at $V_{GS}=1.2$ V the drain current increases only
1.28\% between $V_{DS}=1.0$ and $1.2$ V. The corresponding secant
output conductance is $3.67912\times10^{-5}$ exported current units
per volt. These quantities are descriptive, not model-fit parameters.

As an internal cross-check, the $I_D$--$V_D$ family at
$V_{GS}=V_{DS}=1.2$ V gives $5.8128743\times10^{-4}$
(exported units), consistent to displayed precision with
$5.8128745\times10^{-4}$ in the independently exported $I_D$--$V_G$
transfer sweep at the same bias. The $V_{GS}=0$, $V_{DS}=1.2$ V
operating-point currents likewise agree to displayed precision.
This confirms consistent biasing across the two baseline electrical
simulations.

The baseline electrical simulation requirements of Part 1 are now
fulfilled. Confirmation of the SDevice two-dimensional current
normalization and insertion of the original process screenshots remain
necessary before the report is finalized.
